\begin{topic}[id=bluetooth_mesh,
             difficulty={Mittel},
             type={Protokollvergleich + Systemanalyse},
             prerequisites={Grundlagen Rechnernetze; Basiswissen Funkkommunikation hilfreich},
             practice={optional},
             owner={Dozententeam},
             updated={2026-04-09},
             tags={ Bluetooth Mesh, Provisioning, Relay, Friendship, Security, Models }]{Bluetooth Mesh}

\TopicDescription{Bluetooth Mesh erweitert BLE um vermaschtes Routing – geeignet für Gebäudeautomation und Beleuchtung. Im Thema lernst du Provisioning, Adressierung, Relay/Friendship, Gerätemodelle und Sicherheitsmechanismen. Wir diskutieren Stärken/Schwächen gegenüber Thread/Zigbee und typische Stolpersteine (Energie, Latenz, Skalierung). Ergebnis ist eine fundierte Einschätzung, wo BLE Mesh passt und wie man robuste Netze plant.}

\TopicLeitfragen{\item Welche Grenzen hat BLE Mesh bei Latenz und Skalierung?
\item Wie funktionieren Security und Provisioning zuverlässig?
\item Wann sind Thread/Zigbee die bessere Wahl?
}
\TopicScope{\item Keine Firmware-Portierung.
\item Kein RF-Tuning.
}
\TopicPractice{Optionale Praxisidee: Mini-Mesh (2–3 Nodes) provisionieren und ein einfaches Gerätemodell schalten.}

\TopicKeywords{Bluetooth Mesh, Provisioning, Relay, Friendship, Security, Models}

\nocite{bluetoothMeshSpec,silabsMeshGuide}
\end{topic}
